\documentclass[sigconf,nonacm]{acmart}

\usepackage[brazilian]{babel}
\usepackage{booktabs}
\usepackage{multirow}
\usepackage{subcaption}

\settopmatter{printacmref=false}
\setcopyright{none}
\renewcommand\footnotetextcopyrightpermission[1]{}
\pagestyle{plain}

\title{Comparação de Algoritmos de Filtragem Colaborativa sobre a Base MovieLens 100k}

\author{Mariana Ferreira Ferrarez Ribeiro}
\affiliation{
  \institution{Universidade Federal Rural do Rio de Janeiro}
  \country{Rio de Janeiro, Brasil}
}

\author{Rafael Vieira da Silva}
\affiliation{
  \institution{Universidade Federal Rural do Rio de Janeiro}
  \country{Rio de Janeiro, Brasil}
}

\begin{document}

\begin{abstract}
Este trabalho apresenta a implementação e comparação de oito algoritmos de
filtragem colaborativa, divididos em abordagens baseadas em memória e baseadas
em modelo. Os algoritmos foram avaliados sobre a base MovieLens~100k utilizando
métricas de predição de rating (RMSE, MAE) e de qualidade de recomendação
(Precision@10, Recall@10, NDCG@10, MAP@10). A hiperparametrização foi
conduzida via grid search com validação cruzada de 5 folds, e uma análise
de sensibilidade foi realizada para investigar o impacto dos principais
hiperparâmetros no desempenho. Os resultados mostram que o SVD Regularizado
obteve o melhor desempenho em predição de rating (RMSE = 0,933), enquanto o
BPR apresentou os melhores resultados em métricas de ranking
(Precision@10 = 0,072).
\end{abstract}

\maketitle

%% ============================================================
\section{Introdução}

O crescimento exponencial do comércio eletrônico e das plataformas de conteúdo
digital tem gerado uma sobrecarga de informações que dificulta a tomada de
decisão dos usuários. Nesse cenário, os Sistemas de Recomendação (SR) têm
desempenhado um papel fundamental ao auxiliar os consumidores a encontrar
produtos e serviços relevantes de forma personalizada~\cite{adomavicius2005}.

Grande parte das pesquisas em SR concentra-se na melhoria dos algoritmos de
predição para aprimorar a acurácia das recomendações~\cite{ricci2011}. A
competição Netflix Prize demonstrou que havia oportunidades significativas para
melhorar os algoritmos existentes, impulsionando pesquisas em técnicas como
fatoração matricial, que se tornaram estado da arte~\cite{koren2009}.

O objetivo deste trabalho, desenvolvido no âmbito da disciplina
\emph{Introdução à Teoria de Sistemas de Recomendação}, ministrada pelo
Prof.\ Filipe Braida na UFRRJ, é implementar e comparar diferentes algoritmos
de filtragem colaborativa, incluindo quatro abordagens baseadas em memória e
quatro baseadas em modelo. Além da comparação direta, investigamos a
hiperparametrização de cada algoritmo por meio de grid search com validação
cruzada de 5 folds, e realizamos uma análise de sensibilidade para avaliar o
impacto dos principais hiperparâmetros no desempenho dos modelos. O
código-fonte está disponível publicamente\footnote{https://github.com/ravsil/sistemasDeRecomendacao}.

%% ============================================================
\section{Experimentos}

\subsection{Base de Dados}

Utilizamos a base MovieLens 100k\footnote{https://grouplens.org/datasets/movielens/100k/},
que contém 100.000 avaliações (ratings de 1 a 5) de 943 usuários sobre 1.682
filmes. As avaliações são esparsas --- cada usuário avaliou, em média, cerca de
106 filmes, o que corresponde a uma densidade de aproximadamente 6,3\%.

Para a avaliação final, utilizamos a partição \texttt{ua.base} (90.570
avaliações para treino) e \texttt{ua.test} (9.430 avaliações para teste),
fornecidas pelo próprio dataset. Para a hiperparametrização e análise de
sensibilidade, utilizamos os 5 folds pré-definidos (\texttt{u1}--\texttt{u5}),
cada um com uma divisão 80/20.

\subsection{Algoritmos Implementados}

Todos os algoritmos foram implementados do zero em Python, sem uso de
bibliotecas externas de recomendação. A seguir, descrevemos brevemente cada um.

\subsubsection{Baseados em Memória}

\begin{itemize}
  \item \textbf{Simple Memory}: baseline que combina a média do usuário e a
    média do item com peso $\alpha$, e recomenda itens por popularidade. A
    predição é dada por $\hat{r}_{ui} = \alpha \cdot \bar{r}_u + (1-\alpha) \cdot \bar{r}_i$.
  \item \textbf{User-Based kNN}: calcula a similaridade cosseno entre vetores
    de ratings centrados pela média do usuário. A predição utiliza a média
    ponderada dos desvios dos $k$ vizinhos mais similares que avaliaram o item.
  \item \textbf{Item-Based kNN}: utiliza a similaridade cosseno ajustada
    (adjusted cosine) entre itens, subtraindo a média do usuário antes do
    cálculo. A predição é análoga ao user-based, mas sobre vizinhos de itens.
  \item \textbf{Slope One}: calcula desvios médios entre pares de itens e
    prediz com base nas avaliações do usuário para outros itens, ponderando
    pela contagem de co-ocorrências com suavização (shrinkage):
    $w_{ij} = \frac{c_{ij}}{c_{ij} + \lambda}$, onde $\lambda$ é o fator de shrinkage.
\end{itemize}

\subsubsection{Baseados em Modelo}

\begin{itemize}
  \item \textbf{Simple Model}: modelo de biases que decompõe o rating como
    $\hat{r}_{ui} = \mu + b_u + b_i$, onde $\mu$ é a média global, $b_u$ é o
    bias do usuário e $b_i$ é o bias do item, estimados com regularização $\lambda$.
  \item \textbf{SVD Regularizado}: fatoração matricial com fatores latentes
    treinada via SGD:
    $\hat{r}_{ui} = \mu + b_u + b_i + \mathbf{p}_u^\top \mathbf{q}_i$.
    Os parâmetros são otimizados minimizando o erro quadrático com
    regularização L2 sobre biases e fatores~\cite{koren2009}.
  \item \textbf{BPR (Bayesian Personalized Ranking)}: otimiza um critério
    pairwise para ranking, maximizando a diferença de score entre itens
    positivos (rating $\geq 4$) e negativos amostrados aleatoriamente. Foi
    projetado para feedback implícito, não para predição de rating.
  \item \textbf{Factorization Machines (FM)}: generaliza a fatoração matricial
    usando features one-hot de usuário e item. Modela interações de segunda
    ordem entre features de forma eficiente:
    $\hat{y} = w_0 + \sum_i w_i x_i + \sum_{i<j} \langle \mathbf{v}_i, \mathbf{v}_j \rangle x_i x_j$.
    Treinado via SGD com regularização L2.
\end{itemize}

\subsection{Métricas de Avaliação}

Avaliamos os algoritmos em duas dimensões complementares:

\textbf{Predição de rating}: mede a qualidade da estimativa numérica da nota.
\begin{itemize}
  \item \textbf{RMSE} (Root Mean Squared Error): penaliza erros maiores de
    forma quadrática, sendo mais sensível a outliers.
  \item \textbf{MAE} (Mean Absolute Error): fornece o erro médio absoluto, mais
    interpretável e robusto a outliers que o RMSE.
\end{itemize}

Utilizamos ambas as métricas pois são complementares: o RMSE é o padrão da
literatura para comparação~\cite{koren2009}, enquanto o MAE oferece uma
interpretação mais direta do erro médio.

\textbf{Qualidade de recomendação (top-$K$)}: mede a qualidade da lista
ordenada de recomendações, considerando como relevantes os itens com
rating $\geq 4{,}0$ no conjunto de teste.
\begin{itemize}
  \item \textbf{Precision@$K$}: fração dos itens recomendados que são relevantes.
  \item \textbf{Recall@$K$}: fração dos itens relevantes que foram recomendados.
  \item \textbf{NDCG@$K$}: considera a posição dos acertos, atribuindo maior
    importância a itens relevantes no topo da lista.
  \item \textbf{MAP@$K$}: média da precisão nos pontos de acerto, ponderada
    pela posição.
\end{itemize}

As métricas de ranking são importantes pois, na prática, o usuário visualiza
apenas os primeiros itens recomendados, tornando a ordenação tão relevante
quanto a acurácia pontual.

\subsection{Hiperparametrização}

A hiperparametrização foi realizada via \textbf{grid search com validação
cruzada de 5 folds}, utilizando as partições pré-definidas do MovieLens 100k
(\texttt{u1}--\texttt{u5}). Para cada combinação de hiperparâmetros, o
algoritmo foi treinado e avaliado em cada fold, e o RMSE médio ($\pm$ desvio
padrão) foi registrado.

A Tabela~\ref{tab:hyperparams} apresenta os hiperparâmetros investigados e os
melhores valores encontrados para cada algoritmo.

\begin{table}[h]
\caption{Hiperparâmetros investigados e melhores valores encontrados via
validação cruzada de 5 folds.}
\label{tab:hyperparams}
\scriptsize
\begin{tabular}{@{}p{1.6cm}p{2.0cm}p{2.2cm}p{1.8cm}@{}}
\toprule
\textbf{Algoritmo} & \textbf{Parâmetro} & \textbf{Faixa testada} & \textbf{Melhor} \\
\midrule
Simple Mem. & $\alpha$ (peso) & 0,0 a 1,0 & 0,4 \\
User-Based  & $k$ (vizinhos) & 5 a 80 & 30 \\
Item-Based  & $k$ (vizinhos) & 5 a 80 & 30 \\
Slope One   & $\lambda$ (shrink) & 0 a 150 & 100 \\
\midrule
Simple Mod. & $\lambda$ (reg.) & $10^{-7}$ a 25 & 10 \\
SVD         & $d$, $\eta$, $\lambda$ & 10--50, 0,005--0,01, 0,01--0,05 & 50; 0,01; 0,01 \\
BPR         & $d$, $\eta$, $\lambda$ & 10--50, 0,01--0,05, 0,0025--0,01 & 50; 0,05; 0,0025 \\
FM          & $d$, $\eta$, $\lambda_w$, $\lambda_v$ & 5--20, 0,005--0,01, 0,001--0,01 & 20; 0,01; 0,001; 0,001 \\
\bottomrule
\end{tabular}
\end{table}

\subsection{Resultados}

\subsubsection{Predição de Rating}

A Tabela~\ref{tab:results} apresenta os resultados da avaliação final,
utilizando os melhores hiperparâmetros encontrados.

\begin{table}[h]
\caption{Resultados da avaliação final com hiperparâmetros otimizados
(partição ua.base/ua.test).}
\label{tab:results}
\small
\begin{tabular}{lcccccc}
\toprule
\textbf{Algoritmo} & \textbf{RMSE} & \textbf{MAE} & \textbf{P@10} & \textbf{R@10} & \textbf{NDCG@10} & \textbf{MAP@10} \\
\midrule
Simple Memory & 0,995 & 0,804 & 0,083 & 0,148 & 0,132 & 0,064 \\
User-Based    & 0,959 & 0,750 & 0,002 & 0,003 & 0,002 & 0,000 \\
Item-Based    & 0,956 & 0,747 & 0,010 & 0,019 & 0,013 & 0,005 \\
Slope One     & 0,961 & 0,756 & 0,002 & 0,003 & 0,004 & 0,001 \\
\midrule
Simple Model  & 0,976 & 0,774 & 0,033 & 0,058 & 0,047 & 0,020 \\
\textbf{SVD}  & \textbf{0,933} & \textbf{0,734} & 0,037 & 0,065 & 0,060 & 0,028 \\
BPR           & 2,256 & 1,900 & \textbf{0,072} & \textbf{0,127} & \textbf{0,117} & \textbf{0,057} \\
FM            & 0,959 & 0,759 & 0,032 & 0,055 & 0,039 & 0,014 \\
\bottomrule
\end{tabular}
\end{table}

Em termos de predição de rating, o \textbf{SVD Regularizado} obteve o melhor
desempenho com RMSE de 0,933 e MAE de 0,734. O \textbf{FM} obteve RMSE de
0,959, próximo dos métodos baseados em memória. Os resultados confirmam a
superioridade da fatoração matricial com biases (SVD) para esta tarefa, em
linha com os achados da literatura~\cite{koren2009}.

Entre os métodos baseados em memória, o \textbf{Item-Based kNN} obteve o menor
RMSE (0,956), levemente melhor que o User-Based (0,959) e o Slope One (0,961).
O \textbf{Simple Memory} obteve o pior resultado entre os métodos de memória
(0,995), o que é esperado por ser um baseline simples.

O \textbf{BPR} apresentou RMSE muito alto (2,256) porque foi projetado para
otimizar ranking via feedback implícito, não para predizer ratings. Isso é
esperado e não representa uma falha do algoritmo, mas sim uma inadequação da
métrica para este tipo de modelo.

\subsubsection{Qualidade de Recomendação}

Em métricas de ranking (Precision@10, Recall@10, NDCG@10, MAP@10), o cenário
se inverte: o \textbf{BPR} obteve os melhores resultados em todas as métricas
de ranking (Precision@10 = 0,072, Recall@10 = 0,127), demonstrando que a
otimização pairwise é eficaz para gerar listas de recomendação de qualidade.

O \textbf{Simple Memory} obteve o segundo melhor Precision@10 (0,083), o que
se explica pela sua estratégia de recomendar itens populares --- uma heurística
surpreendentemente competitiva. O \textbf{SVD} ficou em terceiro (0,037),
mostrando que a fatoração matricial também produz boas listas de recomendação.

Notavelmente, os métodos kNN (User-Based e Item-Based) obtiveram métricas de
ranking muito baixas, apesar de boas métricas de predição. Isso sugere que a
acurácia na predição pontual de ratings não se traduz necessariamente em boas
listas de recomendação.

\subsubsection{Visualizações}

As Figuras~\ref{fig:scatter} e~\ref{fig:roc} apresentam, respectivamente, os
scatter plots (rating real vs.\ predito) e as curvas ROC para quatro algoritmos
representativos. A Figura~\ref{fig:confusion} mostra as matrizes de confusão
correspondentes, considerando o limiar de relevância em 4,0.

\begin{figure*}[!t]
\centering
\begin{subfigure}{0.23\textwidth}
  \includegraphics[width=\textwidth]{results/charts/svd/svd_scatter.pdf}
  \caption{SVD}
\end{subfigure}\hfill
\begin{subfigure}{0.23\textwidth}
  \includegraphics[width=\textwidth]{results/charts/fm/fm_scatter.pdf}
  \caption{FM}
\end{subfigure}\hfill
\begin{subfigure}{0.23\textwidth}
  \includegraphics[width=\textwidth]{results/charts/item_based/item_based_scatter.pdf}
  \caption{Item-Based}
\end{subfigure}\hfill
\begin{subfigure}{0.23\textwidth}
  \includegraphics[width=\textwidth]{results/charts/bpr/bpr_scatter.pdf}
  \caption{BPR}
\end{subfigure}
\caption{Scatter plots: rating real vs.\ predito.}
\label{fig:scatter}
\end{figure*}

\begin{figure*}[!t]
\centering
\begin{subfigure}{0.23\textwidth}
  \includegraphics[width=\textwidth]{results/charts/svd/svd_roc.pdf}
  \caption{SVD}
\end{subfigure}\hfill
\begin{subfigure}{0.23\textwidth}
  \includegraphics[width=\textwidth]{results/charts/fm/fm_roc.pdf}
  \caption{FM}
\end{subfigure}\hfill
\begin{subfigure}{0.23\textwidth}
  \includegraphics[width=\textwidth]{results/charts/item_based/item_based_roc.pdf}
  \caption{Item-Based}
\end{subfigure}\hfill
\begin{subfigure}{0.23\textwidth}
  \includegraphics[width=\textwidth]{results/charts/bpr/bpr_roc.pdf}
  \caption{BPR}
\end{subfigure}
\caption{Curvas ROC (limiar = 4,0).}
\label{fig:roc}
\end{figure*}

\begin{figure*}[!t]
\centering
\begin{subfigure}{0.23\textwidth}
  \includegraphics[width=\textwidth]{results/charts/svd/svd_confusion.pdf}
  \caption{SVD}
\end{subfigure}\hfill
\begin{subfigure}{0.23\textwidth}
  \includegraphics[width=\textwidth]{results/charts/fm/fm_confusion.pdf}
  \caption{FM}
\end{subfigure}\hfill
\begin{subfigure}{0.23\textwidth}
  \includegraphics[width=\textwidth]{results/charts/item_based/item_based_confusion.pdf}
  \caption{Item-Based}
\end{subfigure}\hfill
\begin{subfigure}{0.23\textwidth}
  \includegraphics[width=\textwidth]{results/charts/bpr/bpr_confusion.pdf}
  \caption{BPR}
\end{subfigure}
\caption{Matrizes de confusão (limiares = 4,0).}
\label{fig:confusion}
\end{figure*}

\subsubsection{Análise de Sensibilidade}

A análise de sensibilidade investigou como a variação de um hiperparâmetro-chave
impacta o desempenho (RMSE) de cada algoritmo, mantendo os demais fixos nos
melhores valores. Os principais achados foram:

\textbf{Número de vizinhos $k$ (User-Based e Item-Based)}: ambos os algoritmos
apresentaram uma curva em forma de U, com RMSE diminuindo rapidamente de $k=5$
até $k=20$--$30$, estabilizando e depois aumentando levemente para $k>50$. O
ponto ótimo $k=30$ indica um equilíbrio entre incluir vizinhos informados e
evitar ruído de vizinhos pouco similares.

\textbf{Número de fatores latentes (SVD, BPR, FM)}: para os três modelos, mais
fatores latentes resultaram em melhor desempenho, com ganhos decrescentes. O SVD
passou de 0,928 ($d=5$) para 0,911 ($d=100$), uma melhoria de 1,8\%. O FM
apresentou comportamento similar, mas com ganho menos monotônico.

\textbf{Shrinkage $\lambda$ do Slope One}: o RMSE diminuiu monotonicamente com
o aumento de $\lambda$ (de 0,948 sem shrinkage para 0,943 com $\lambda=150$),
indicando que a suavização por contagem é benéfica e que pares com poucos
co-raters devem ser penalizados.

\textbf{Peso $\alpha$ do Simple Memory}: a curva mostrou um formato de U
simétrico com mínimo em $\alpha=0{,}4$, indicando que dar peso ligeiramente
maior à média do item ($60\%$) do que à do usuário ($40\%$) é ótimo. Isso sugere
que a qualidade intrínseca do item é um preditor ligeiramente mais informativo
que o padrão geral do usuário.

\textbf{Regularização do Simple Model}: o RMSE decresceu até $\lambda=10$ e
depois voltou a subir para $\lambda=25$, mostrando uma curva em U típica do
trade-off entre viés e variância.

As Figuras~\ref{fig:sens_memory} a~\ref{fig:sens_model} ilustram graficamente
essas tendências.

\begin{figure*}[!t]
\centering
\begin{subfigure}{0.23\textwidth}
  \includegraphics[width=\textwidth]{results/sensitivity/simple_memory_alpha.pdf}
  \caption{Simple Mem.}
\end{subfigure}\hfill
\begin{subfigure}{0.23\textwidth}
  \includegraphics[width=\textwidth]{results/sensitivity/user_based_k.pdf}
  \caption{User-Based}
\end{subfigure}\hfill
\begin{subfigure}{0.23\textwidth}
  \includegraphics[width=\textwidth]{results/sensitivity/item_based_k.pdf}
  \caption{Item-Based}
\end{subfigure}\hfill
\begin{subfigure}{0.23\textwidth}
  \includegraphics[width=\textwidth]{results/sensitivity/slope_one_shrink.pdf}
  \caption{Slope One}
\end{subfigure}
\caption{Sensibilidade -- algoritmos baseados em memória.}
\label{fig:sens_memory}
\end{figure*}

\begin{figure*}[!t]
\centering
\begin{subfigure}{0.23\textwidth}
  \includegraphics[width=\textwidth]{results/sensitivity/simple_model_reg.pdf}
  \caption{Simple Model}
\end{subfigure}\hfill
\begin{subfigure}{0.23\textwidth}
  \includegraphics[width=\textwidth]{results/sensitivity/svd_n_factors.pdf}
  \caption{SVD}
\end{subfigure}\hfill
\begin{subfigure}{0.23\textwidth}
  \includegraphics[width=\textwidth]{results/sensitivity/bpr_n_factors.pdf}
  \caption{BPR}
\end{subfigure}\hfill
\begin{subfigure}{0.23\textwidth}
  \includegraphics[width=\textwidth]{results/sensitivity/fm_n_factors.pdf}
  \caption{FM}
\end{subfigure}
\caption{Sensibilidade -- algoritmos baseados em modelo.}
\label{fig:sens_model}
\end{figure*}

%% ============================================================
\section{Conclusão}

Este trabalho implementou e comparou oito algoritmos de filtragem colaborativa,
divididos entre abordagens baseadas em memória e baseadas em modelo, sobre a
base MovieLens 100k.

Os principais achados foram: (1)~o \textbf{SVD Regularizado} obteve o melhor
desempenho em predição de rating (RMSE = 0,933), confirmando a eficácia da
fatoração matricial; (2)~o \textbf{BPR} superou todos os outros em métricas de
ranking, demonstrando a importância de otimizar diretamente para a tarefa
desejada; (3)~a acurácia na predição de rating não implica necessariamente em
boas recomendações, como evidenciado pelos métodos kNN; e (4)~a
hiperparametrização produziu melhorias consistentes, com o SVD reduzindo seu
RMSE de 0,958 (default) para 0,933 (otimizado), uma melhoria de 2,6\%.

A análise de sensibilidade revelou padrões interpretáveis: o número ideal de
vizinhos para métodos kNN é finito ($k \approx 30$), enquanto modelos de
fatoração se beneficiam de mais fatores latentes com retornos decrescentes.

Como sugestões para trabalhos futuros, destacamos: (a)~explorar técnicas de
ensemble combinando modelos de predição e de ranking; (b)~avaliar os algoritmos
em bases maiores (MovieLens 1M, CiaoDVD) para testar a escalabilidade;
(c)~incorporar informações contextuais e de conteúdo para melhorar as
recomendações cold-start; e (d)~investigar técnicas de preprocessamento como
normalização de ratings e remoção de outliers.

\vfill
%% ============================================================
\bibliographystyle{ACM-Reference-Format}
\begin{thebibliography}{3}

\bibitem{adomavicius2005}
G.~Adomavicius and A.~Tuzhilin.
\newblock Toward the next generation of recommender systems: a survey of the
  state-of-the-art and possible extensions.
\newblock \emph{IEEE Transactions on Knowledge and Data Engineering},
  17(6):734--749, jun 2005. ISSN 1041-4347. doi: 10.1109/TKDE.2005.99.

\bibitem{koren2009}
Y.~Koren, R.~Bell, and C.~Volinsky.
\newblock Matrix factorization techniques for recommender systems.
\newblock \emph{Computer}, 42(8):30--37, 2009.

\bibitem{ricci2011}
F.~Ricci, L.~Rokach, B.~Shapira, and P.~B. Kantor.
\newblock \emph{Recommender Systems Handbook}.
\newblock Springer, 2011. doi: 10.1007/978-0-387-85820-3.

\end{thebibliography}

\end{document}
